\chapter{Tutorials}
\label{ch:tutorials}

This chapter provides step-by-step tutorials for common tasks with
Veilbox, from first build to advanced container operations.

\section{Tutorial 1: First Build}

\textbf{Goal}: Build the Veilbox disk image from scratch.

\begin{enumerate}[nosep]
  \item Clone the repository:
\begin{lstlisting}[language=bash]
git clone https://github.com/Shreyas0047/veilbox.git
cd veilbox
\end{lstlisting}
  \item Check prerequisites:
\begin{lstlisting}[language=bash]
gcc --version          # Need 14+
make --version         # Need 4.4+
qemu-system-x86_64 --version  # Need 7.0+
\end{lstlisting}
  \item Run the build:
\begin{lstlisting}[language=bash]
./build.sh
# Full build: ~5 minutes
# Output: output/veilbox.raw, output/veilbox.vdi, output/vmlinuz
\end{lstlisting}
  \item Boot in QEMU:
\begin{lstlisting}[language=bash]
MEM=2G ./test.sh
\end{lstlisting}
  \item When you see the login prompt, log in:
\begin{lstlisting}
veilbox login: root
Password: veiladmin
\end{lstlisting}
\end{enumerate}

\section{Tutorial 2: Quick Rebuild}

\textbf{Goal}: Modify a configuration file and rebuild quickly.

\begin{enumerate}[nosep]
  \item Edit the rcS script:
\begin{lstlisting}[language=bash]
vim rootfs/etc/init.d/rcS
\end{lstlisting}
  \item Rebuild initramfs and kernel:
\begin{lstlisting}[language=bash]
# Regenerate initramfs file list
./build.sh --stage initramfs

# Rebuild kernel with updated initramfs
./build.sh --stage kernel

# Rebuild disk image
./build.sh --stage disk
\end{lstlisting}
  \item Boot with updated image:
\begin{lstlisting}[language=bash]
MEM=2G ./test.sh
\end{lstlisting}
\end{enumerate}

\section{Tutorial 3: SSH Access}

\textbf{Goal}: Connect to Veilbox via SSH from the host.

\begin{enumerate}[nosep]
  \item Boot Veilbox with SSH port forwarding:
\begin{lstlisting}[language=bash]
MEM=2G ./test.sh
\end{lstlisting}
  \item From another terminal, connect:
\begin{lstlisting}[language=bash]
ssh -i output/ssh-test-key root@localhost -p 2222
\end{lstlisting}
  \item Or with password:
\begin{lstlisting}[language=bash]
ssh root@localhost -p 2222
# Password: veiladmin
\end{lstlisting}
\end{enumerate}

\section{Tutorial 4: Run a Container}

\textbf{Goal}: Pull and run an nginx container.

\begin{enumerate}[nosep]
  \item Log in to Veilbox (console or SSH).
  \item Check containerd is running:
\begin{lstlisting}
ps | grep containerd
\end{lstlisting}
  \item Pull an image:
\begin{lstlisting}
nerdctl pull nginx:alpine
nerdctl images
\end{lstlisting}
  \item Run the container:
\begin{lstlisting}
nerdctl run -d --name web -p 80:80 nginx:alpine
nerdctl ps
\end{lstlisting}
  \item Test from Veilbox:
\begin{lstlisting}
wget -q -O- http://localhost
\end{lstlisting}
\end{enumerate}

\section{Tutorial 5: Persistent Storage}

\textbf{Goal}: Create a container with persistent data.

\begin{enumerate}[nosep]
  \item Verify state partition is mounted:
\begin{lstlisting}
df -h /mnt/state
\end{lstlisting}
  \item Create a data directory:
\begin{lstlisting}
mkdir -p /mnt/state/data
\end{lstlisting}
  \item Write data:
\begin{lstlisting}
echo "Hello from state" > /mnt/state/data/hello.txt
\end{lstlisting}
  \item Run a container with the data mounted:
\begin{lstlisting}
nerdctl run --rm -v /mnt/state/data:/data alpine \
    cat /data/hello.txt
\end{lstlisting}
  \item Reboot and verify persistence:
\begin{lstlisting}
MEM=2G ./test.sh --keep-state
# After login:
cat /mnt/state/data/hello.txt
\end{lstlisting}
\end{enumerate}


\section{Tutorial 6: Custom Kernel Configuration}

\textbf{Goal}: Add a new kernel feature and rebuild.

\begin{enumerate}[nosep]
  \item Enter kernel configuration:
\begin{lstlisting}[language=bash]
cd linux-7.1.0-rc6
make menuconfig
\end{lstlisting}
  \item Navigate to the desired feature. For example, enable BPF:
        \texttt{General setup} $\rightarrow$ \texttt{BPF subsystem}.
  \item Save the configuration:
\begin{lstlisting}[language=bash]
# Save as custom fragment
grep -E "^CONFIG_BPF" .config > ../kernel/configs/bpf.config
\end{lstlisting}
  \item Merge into main config:
\begin{lstlisting}[language=bash]
./scripts/kconfig/merge_config.sh \
    -O .config \
    .config \
    ../kernel/configs/bpf.config
make olddefconfig
\end{lstlisting}
  \item Rebuild:
\begin{lstlisting}[language=bash]
make -j$(nproc) bzImage
cp arch/x86/boot/bzImage ../output/vmlinuz
\end{lstlisting}
  \item Test the new feature:
\begin{lstlisting}
# If BPF was added:
mount -t bpf none /sys/fs/bpf
bpftool prog list
\end{lstlisting}
\end{enumerate}

\section{Tutorial 7: Debugging Build Failures}

\textbf{Goal}: Diagnose and fix a common kernel build error.

\begin{enumerate}[nosep]
  \item Run the build:
\begin{lstlisting}[language=bash]
make -j$(nproc) bzImage 2>&1 | tee build.log
\end{lstlisting}
  \item Search for errors:
\begin{lstlisting}[language=bash]
grep -i error build.log | head -20
\end{lstlisting}
  \item Common fixes:
\begin{itemize}[nosep]
  \item \texttt{implicit declaration}: Missing header or wrong GCC version.
  \item \texttt{undefined reference}: Missing library or wrong config.
  \item \texttt{too many sections}: Disable CONFIG\_DEBUG\_INFO.
  \item \texttt{disk full}: Free space in /tmp or build directory.
\end{itemize}
  \item If the error is config-related, use:
\begin{lstlisting}[language=bash]
make olddefconfig    # Fix dependency issues
make menuconfig      # Manually fix options
\end{lstlisting}
  \item Clean and retry:
\begin{lstlisting}[language=bash]
make clean
make -j$(nproc) bzImage
\end{lstlisting}
\end{enumerate}

\section{Tutorial 8: Booting from USB}

\textbf{Goal}: Write Veilbox to a USB drive for bare-metal boot.

\begin{enumerate}[nosep]
  \item Insert a USB drive (8 GB or larger).
  \item Identify the device:
\begin{lstlisting}[language=bash]
lsblk
# Look for your USB device, e.g., /dev/sdX
\end{lstlisting}
  \item Write the disk image:
\begin{lstlisting}[language=bash]
sudo dd if=output/veilbox.raw of=/dev/sdX bs=4M status=progress
sync
\end{lstlisting}
  \item Boot from the USB drive on the target machine (configure BIOS
        boot order to boot from USB).
  \item Log in with \texttt{root} / \texttt{veiladmin}.
\end{enumerate}

\section{Tutorial 9: Creating Custom Container Images}

\textbf{Goal}: Build a custom container image with nerdctl.

\begin{enumerate}[nosep]
  \item Create a Dockerfile:
\begin{lstlisting}
mkdir /mnt/state/myapp
cat > /mnt/state/myapp/Dockerfile << 'EOF'
FROM alpine:latest
RUN apk add --no-cache python3 py3-pip
COPY app.py /app.py
CMD ["python3", "/app.py"]
EOF
\end{lstlisting}
  \item Build the image:
\begin{lstlisting}
nerdctl build -t myapp:latest /mnt/state/myapp
\end{lstlisting}
  \item Run the image:
\begin{lstlisting}
nerdctl run -d --name myapp -p 5000:5000 myapp:latest
\end{lstlisting}
   \item Push to a registry (optional):
\begin{lstlisting}
nerdctl tag myapp:latest registry.example.com/myapp:latest
nerdctl push registry.example.com/myapp:latest
\end{lstlisting}
\end{enumerate}

\section{Tutorial 10: Host Your Own Website}

\textbf{Goal}: Containerize and host your own website on Veilbox.

\begin{enumerate}[nosep]
  \item Transfer your website folder to the VM:
\begin{lstlisting}[language=bash]
# From host
scp -P 2222 -r /path/to/website root@localhost:/mnt/state/site
\end{lstlisting}
  \item Pull a web server image:
\begin{lstlisting}
nerdctl pull nginx:alpine
\end{lstlisting}
  \item Serve your site with nginx (volume mount):
\begin{lstlisting}
nerdctl run -d --name web \
    --network host \
    -v /mnt/state/site:/usr/share/nginx/html:ro \
    nginx:alpine
\end{lstlisting}
  \item View from host at \url{http://localhost:8080/}.
  \item For a Node.js app with a backend:
\begin{lstlisting}
nerdctl run -d --name myapp \
    --network host \
    -e PORT=8080 \
    -v /mnt/state/myapp:/app:ro \
    -w /app \
    node:26-slim \
    node server.js
\end{lstlisting}
  \item Update live by re-transferring files over SCP --- no restart needed.
\end{enumerate}

\textbf{Note}: Use \texttt{--network host} instead of port mapping (\texttt{-p 8080:3000})
because QEMU user-mode networking does not support kernel bridge operations
(docker-style port forwarding). The \texttt{-e PORT=8080} environment variable
makes the app listen on the port that is already forwarded by \texttt{test.sh}
to the host.
